% Patch record for C-61 -- debate inconsistencies S-02, S-05, S-06, S-07, checked against code/data.
% Applied 2026-10-29 (Sebas). Source: intake/pending/second_pass_simplification_debate.md.
% Evidence (read-only, sammcar/PETER_SIMULATION@main via gh api; experiments/simulation/):
%  S-02 noise phrases kept. Author confirmed (2026-10-29) the noise trials were run; setup inferred from
%       familia_c1/c2 suite_execution_manifest.json (noise_level_idx cycles 0-4, 3 trials per level) and
%       red_neuronal.py NoiseLevel = [0, 0.05, 0.10, 0.20, 0.30] (proportional Gaussian on IMU, LiDAR,
%       camera position/area). Setup sentence added; "unperturbed" in sssec:noise_robustness_extra fixed.
%  S-05 largest reported means: Roll 1.59 deg (rugged, fig:terrain_adaptability), Pitch 2.29 deg (slope).
%  S-06 red_neuronal.py: z[1] = f(pitch - Upitch), z[16] += 5*z[1] (z16 -> mode 'H' = differential rolling);
%       familia_c2_pendiente/*/metrics_raw.csv mode column ends in H; NEU_IMU2.png X16 replaces X15 at ~38 s;
%       NEU_IMU2_real.png X16 (X14 never active). Table Locom_NN thresholds left as published.
%  S-07 table:HardwareParams: camera is a Logitech C270 (no depth); obstacles are LiDAR-detected.
% results.tex l.945 also carries one C-62 sentence (the step-response sentence); see c-62 record.
% Open for authors: tab:terrain_sequence Frame 1 mode on the ramp.
% validate_tex.sh OK; check_roundtrip.sh OK; compile_pdf.sh OK (0 undefined references).
%
% --- a/sections/conclusions.tex
% +++ b/sections/conclusions.tex
% @@ -10 +10 @@ In scenarios presenting simultaneous appetitive and aversive stimuli, the basal-
% -On irregular surfaces and steep slopes, the system exhibited an adaptive response grounded in the direct integration of inertial data into the locomotion control module. Continuous monitoring of roll, pitch, and the standard deviation of vertical acceleration ($\sigma_{az}$) enabled early detection of instability conditions. When the predefined thresholds for each metric were exceeded, the neural architecture immediately triggered a transition from H-mode to quadrupedal mode, characterized by a gradual activation pattern in the modulator neurons ($X_{0}$–$X_{2}$) and a controlled inhibition of the previous gait. This transition not only maintained postural stability but also significantly reduced lateral and longitudinal oscillation amplitudes, minimizing the risk of loss of balance. The relationship between IMU signal intensity and postural-response activation confirms that combining kinematic and dynamic criteria is an effective strategy for locomotor adaptation in adverse terrains, ensuring both structural integrity of the robot and continuity of the mission. Likewise, it is noted that the first version of the gait-decision network implemented in simulation produces variant transient states during transitions to and from the quadrupedal mode; therefore, exploring additional IMU-derived variables or signal-processing approaches could yield more robust terrain-state inference. 
% +On irregular surfaces and steep slopes, the system exhibited an adaptive response grounded in the direct integration of inertial data into the locomotion control module. Continuous monitoring of roll, pitch, and the standard deviation of vertical acceleration ($\sigma_{az}$) enabled early detection of instability conditions. \revblue{When a threshold was exceeded, the neural architecture switched gait: high $\sigma_{az}$ on rugged terrain engaged the quadrupedal mode, while sustained tilt on the slope engaged differential rolling. Each switch showed a gradual activation of the modulator neurons ($X_{0}$–$X_{2}$) and a controlled inhibition of the previous gait. These transitions} not only maintained postural stability but also significantly reduced lateral and longitudinal oscillation amplitudes, minimizing the risk of loss of balance. The relationship between IMU signal intensity and postural-response activation confirms that combining kinematic and dynamic criteria is an effective strategy for locomotor adaptation in adverse terrains, ensuring both structural integrity of the robot and continuity of the mission. Likewise, it is noted that the first version of the gait-decision network implemented in simulation produces variant transient states during transitions to and from the quadrupedal mode; therefore, exploring additional IMU-derived variables or signal-processing approaches could yield more robust terrain-state inference. 
% --- a/sections/methodology.tex
% +++ b/sections/methodology.tex
% @@ -323 +323 @@ where $x$ is the net input to the neuron; $x_+ = \max(x,0)$ is the rectified inp
% -\tau \frac{dz_1}{dt} &= -z_1 +  f_{Imu}( Roll - U_{Roll})
% +\tau \frac{dz_1}{dt} &= -z_1 +  f_{Imu}( \revblue{Pitch - U_{Pitch}})
% @@ -325 +325 @@ where $x$ is the net input to the neuron; $x_+ = \max(x,0)$ is the rectified inp
% -\tau \frac{dz_2}{dt} &= -z_2 +  f_{Imu}( Pitch - U_{Pitch})
% +\tau \frac{dz_2}{dt} &= -z_2 +  f_{Imu}( \revblue{Roll - U_{Roll}})
% --- a/sections/results.tex
% +++ b/sections/results.tex
% @@ -152 +152 @@ Complex Navigation   & 100.0\% & 24.96 $\pm$ 10.36 & 1.03 $\pm$ 0.35 & 1.01 $\pm
% -Two further scenarios --- the Rugged Terrain course and the Inclined Slope, already introduced in Section~\ref{sssec:variable_topo_quant} --- were re-evaluated under the same combined IMU-drift/LiDAR-noise/illumination-distortion sweep as the four scenarios above (Section~\ref{ssec:NoiseRobustnessEval}). This is a separate battery from the unperturbed Rugged Terrain/Inclined Slope trials analyzed in Section~\ref{sssec:variable_topo_quant} and is not directly comparable to it; the two batteries differ in perturbation condition, not just sample size.
% +Two further scenarios --- the Rugged Terrain course and the Inclined Slope, already introduced in Section~\ref{sssec:variable_topo_quant} --- were re-evaluated under the same combined IMU-drift/LiDAR-noise/illumination-distortion sweep as the four scenarios above (Section~\ref{ssec:NoiseRobustnessEval}). This is a separate battery from the Rugged Terrain/Inclined Slope trials analyzed in Section~\ref{sssec:variable_topo_quant}\revblue{, which carried only proportional Gaussian sensor noise,} and is not directly comparable to it; the two batteries differ in perturbation condition, not just sample size.
% @@ -240 +240 @@ In the first simulation experiment shown in Figure~\ref{fig:RES_IMU}, the platfo
% -For the second stability test in simulation, the neural activities of the gait-mode control structure were recorded (Figure~\ref{fig:NEU_IMU2}b). At $t=41$~seconds, neuron $X_{1}$---associated with detecting an inclination greater than five degrees---excites unit $X_{15}$, issuing the command to transform into quadrupedal mode. The transition during locomotion is illustrated in Figure~\ref{fig:RES_IMU2}, with snapshots of the platform posture taken every 12 seconds from a lateral perspective. It is important to highlight the neural responses shown in Figure~\ref{fig:NEU_IMU2}a, where the persistent activation of neuron $X_{4}$ produces sustained activity in neurons $X_{11}$ and $X_{12}$, which in turn drive neuron $X_{13}$ and ultimately generate the forward escape behavior. In Figure~\ref{fig:GRAPH_IMU2}, the data recorded by the inertial sensor are shown.
% +For the second stability test in simulation, the neural activities of the gait-mode control structure were recorded (Figure~\ref{fig:NEU_IMU2}b). \revblue{From $t \approx 38$~s, neuron $X_{1}$, which responds to pitch above the $U_{Pitch}$ threshold, excites unit $X_{16}$; $X_{16}$ suppresses $X_{15}$ and commands the switch from quadrupedal to differential rolling mode.} The transition during locomotion is illustrated in Figure~\ref{fig:RES_IMU2}, with snapshots of the platform posture taken every 12 seconds from a lateral perspective. It is important to highlight the neural responses shown in Figure~\ref{fig:NEU_IMU2}a, where the persistent activation of neuron $X_{4}$ produces sustained activity in neurons $X_{11}$ and $X_{12}$, which in turn drive neuron $X_{13}$ and ultimately generate the forward escape behavior. In Figure~\ref{fig:GRAPH_IMU2}, the data recorded by the inertial sensor are shown.
% @@ -320 +320 @@ For the second stability test in simulation, the neural activities of the gait-m
% -To quantify the computational efficiency and temporal determinism of the proposed bio-inspired control architecture under stochastic conditions, the neural response latency ($T_{\text{response}}$) and the Empirical Cumulative Distribution Function (ECDF) of the mechanical state-switching delay ($T_{\text{switch}}$) were analyzed across 15 randomized trials for both experimental terrains, the latter illustrated in Figure~\ref{fig:decision_delay}. \revblue{The neural latency was} $T_{\text{response}} = 0.0\,\text{s}$ for both terrain types, showing that the threshold-based basal ganglia selection logic filtered out the injected sensor noise and dispatched the correct locomotion command within the same simulation frame as the environmental disturbance. The mechanical transition latency was likewise sub-millisecond in the vast majority of trials, but not fully invariant across terrain types. On rugged terrain, $T_{\text{switch}}$ fell within $0.2$--$0.7\,\text{ms}$ for 13 of 15 trials (a single discrete update cycle of the ROS~2 Control loop operating at $3.3\,\text{kHz}$), while the remaining 2 trials ($13.3\%$) exhibited a long tail out to $\approx 0.6\,\text{s}$, indicating that occasional trials require additional update cycles to complete the locomotion-matrix hot-swap under this terrain's contact conditions. On the inclined slope, $T_{\text{switch}}$ spread across $0.9$--$1.8\,\text{ms}$ in a five-to-six-step staircase rather than collapsing onto a single value, reflecting genuine trial-to-trial dispersion of a comparable order of magnitude rather than a perfectly deterministic delay. \revblue{Neural decision dispatch thus remained deterministic across all 30 trials, while the residual mechanical variability, including the rugged-terrain long tail, is a bounded, terrain-dependent effect rather than chattering.}
% +To quantify the computational efficiency and temporal determinism of the proposed bio-inspired control architecture under stochastic conditions, the neural response latency ($T_{\text{response}}$) and the Empirical Cumulative Distribution Function (ECDF) of the mechanical state-switching delay ($T_{\text{switch}}$) were analyzed across 15 randomized trials for both experimental terrains, the latter illustrated in Figure~\ref{fig:decision_delay}. \revblue{In these trials, zero-mean Gaussian noise proportional to the measured value was added to the IMU, LiDAR, and camera inputs at five levels ($\sigma = 0$, 5, 10, 20, and 30\%), with three trials per level on each terrain.} \revblue{The neural latency was} $T_{\text{response}} = 0.0\,\text{s}$ for both terrain types, showing that the threshold-based basal ganglia selection logic filtered out the injected sensor noise and dispatched the correct locomotion command within the same simulation frame as the environmental disturbance. The mechanical transition latency was likewise sub-millisecond in the vast majority of trials, but not fully invariant across terrain types. On rugged terrain, $T_{\text{switch}}$ fell within $0.2$--$0.7\,\text{ms}$ for 13 of 15 trials (a single discrete update cycle of the ROS~2 Control loop operating at $3.3\,\text{kHz}$), while the remaining 2 trials ($13.3\%$) exhibited a long tail out to $\approx 0.6\,\text{s}$, indicating that occasional trials require additional update cycles to complete the locomotion-matrix hot-swap under this terrain's contact conditions. On the inclined slope, $T_{\text{switch}}$ spread across $0.9$--$1.8\,\text{ms}$ in a five-to-six-step staircase rather than collapsing onto a single value, reflecting genuine trial-to-trial dispersion of a comparable order of magnitude rather than a perfectly deterministic delay. \revblue{Neural decision dispatch thus remained deterministic across all 30 trials, while the residual mechanical variability, including the rugged-terrain long tail, is a bounded, terrain-dependent effect rather than chattering.}
% @@ -466,2 +466,3 @@ characterize the locomotion phases immediately preceding and
% -following each transition; these values remain below $1.36^\circ$
% -and $6.95^\circ$ respectively across all tested conditions,
% +following each transition; \revblue{the mean Roll RMS stays at or
% +below $1.59^\circ$ and the mean Pitch RMS at or below $2.29^\circ$
% +across these suites and both variable-topography terrains},
% @@ -757 +758 @@ In the real implementation (Figure~\ref{fig:RES_IMU_real}), images were captured
% -A similar behavior is observed in the experiment with the physical robot (Figure~\ref{fig:RES_IMU2_real}). However, in the real robot's gait-decision network, the inclined terrain is detected before second 46 (Figure~\ref{fig:NEU_IMU2_real}, Locomotion Module panel); consequently, a transient response emerges, increasing the activity of neuron $X_{14}$, which commands the robot to climb the ramp using the differential rolling mode; this mode remains active even after clearing the slope. \revblue{MLP outputs and stagnation state for both scenarios are shown in Figures~\ref{fig:GRAPH_IMU_real} and~\ref{fig:GRAPH_IMU2_real}.}
% +A similar behavior is observed in the experiment with the physical robot (Figure~\ref{fig:RES_IMU2_real}). However, in the real robot's gait-decision network, the inclined terrain is detected before second 46 (Figure~\ref{fig:NEU_IMU2_real}, Locomotion Module panel); consequently, a transient response emerges, increasing the activity of neuron \revblue{$X_{16}$}, which commands the robot to climb the ramp using the differential rolling mode; this mode remains active even after clearing the slope. \revblue{MLP outputs and stagnation state for both scenarios are shown in Figures~\ref{fig:GRAPH_IMU_real} and~\ref{fig:GRAPH_IMU2_real}.}
% @@ -898 +899 @@ To further consolidate the physical validation under complex, multi-section topo
% -\textbf{Appetitive (see Figure~\ref{fig:appetitive_real_compound}):} The extended switching-delay window does not stem from computational overhead within the basal ganglia network---which operates at the microsecond scale---but from the electro-mechanical trajectory smoothing and landing interpolation enforced by the low-level controllers during target approach, to prevent peak-torque damage on the 3D-printed actuators. The prolonged pre-lock exploration phase reflects real-world camera field-of-view and illumination constraints; the sustained, non-zero firing variance after mission completion reflects that the stopping neuron ($X_{17}$) freezes the physical stance while upstream sensory pathways continue registering the stimulus, rather than a full sensory shutdown. The cyclic Roll/Pitch oscillations stem from body-shaking during the rhythmic weight-shifting quadrupedal gait over uncalibrated floor textures, flattening once the platform converges to a stable equilibrium.
% +\textbf{Appetitive (see Figure~\ref{fig:appetitive_real_compound}):} \revblue{The extended switching-delay window comes mainly from the electro-mechanical trajectory smoothing and landing interpolation enforced by the low-level controllers during target approach, which prevent peak-torque damage to the servo actuators.} The prolonged pre-lock exploration phase reflects real-world camera field-of-view and illumination constraints; the sustained, non-zero firing variance after mission completion reflects that the stopping neuron ($X_{17}$) freezes the physical stance while upstream sensory pathways continue registering the stimulus, rather than a full sensory shutdown. The cyclic Roll/Pitch oscillations stem from body-shaking during the rhythmic weight-shifting quadrupedal gait over uncalibrated floor textures, flattening once the platform converges to a stable equilibrium.
% @@ -944 +945 @@ To further consolidate the physical validation under complex, multi-section topo
% -\textbf{Complex (see Figure~\ref{fig:complex_real_compound}):} \revblue{Across successive back-to-back transformations, the deterministic step response shows the ROS~2 Control architecture managing electromechanical inertia and trajectory blending, protecting the joints from cumulative torque stress.} The three-phase firing-variance profile maps the multi-stage decision process: a transient hyper-excitation peak during static obstacle avoidance, driven by the physical barrier's high spatial density flooding the camera sensor; a secondary activation wave capturing the emergency flight reflex to the aversive cue; and a low-energy steady-state as the robot tracks the appetitive target to mission completion. The chassis attitude telemetry mirrors this sequence, with asymmetric flanking maneuvers exciting the transversal axis and rapid backward accelerations peaking the longitudinal tilt, before both metrics drop and stabilize during the final convergence phase, indicating the wheel-driven configuration smooths parasitic oscillations during target approach.
% +\textbf{Complex (see Figure~\ref{fig:complex_real_compound}):} \revblue{The step response stays deterministic across back-to-back transformations: the ROS~2 Control layer absorbs the electromechanical inertia and blends the trajectories, so torque stress does not accumulate in the joints.} The three-phase firing-variance profile maps the multi-stage decision process: a transient hyper-excitation peak during static obstacle avoidance, \revblue{driven by the dense LiDAR returns from the physical barrier}; a secondary activation wave capturing the emergency flight reflex to the aversive cue; and a low-energy steady-state as the robot tracks the appetitive target to mission completion. The chassis attitude telemetry mirrors this sequence, with asymmetric flanking maneuvers exciting the transversal axis and rapid backward accelerations peaking the longitudinal tilt, before both metrics drop and stabilize during the final convergence phase, indicating the wheel-driven configuration smooths parasitic oscillations during target approach.
% @@ -967 +968 @@ To further consolidate the physical validation under complex, multi-section topo
% -\textbf{Obstacle (see Figure~\ref{fig:obstacle_real_compound}):} \revblue{As in the Aversive scenario, the step response is immediate and consistent across all three replicates: geometric blockages are handled as emergency reflexes without trajectory ramp filtering.} The massive transient spike in firing variance---an order of magnitude above the color-cue scenarios---reflects the volumetric obstacle flooding the camera sensor with spatial and depth data, triggering a hyper-excitation wave that drives aggressive lateral flanking, collapsing back to baseline once the chassis clears the obstacle's spatial boundary. The rhythmic Roll expansion with coupled Pitch variation reflects the weight-shifting needed for a tight lateral bypass, flattening once safety is secured.
% +\textbf{Obstacle (see Figure~\ref{fig:obstacle_real_compound}):} \revblue{As in the Aversive scenario, the step response is immediate and consistent across all three replicates: geometric blockages are handled as emergency reflexes without trajectory ramp filtering.} The massive transient spike in firing variance---an order of magnitude above the color-cue scenarios---\revblue{reflects the obstacle entering the LiDAR detection range and activating its obstacle-encoding ring}, triggering a hyper-excitation wave that drives aggressive lateral flanking, collapsing back to baseline once the chassis clears the obstacle's spatial boundary. The rhythmic Roll expansion with coupled Pitch variation reflects the weight-shifting needed for a tight lateral bypass, flattening once safety is secured.
